\section{A short integral basis for the period lattice}
\label{sec:period-basis}

We prove the numerical lattice estimate used in the jet comparison.
Let \(E\) be the CM host field, \(d=[E:\mathbb Q]\),
\(\Delta=|\operatorname{disc}E|\), and \(R=\Delta^{2/d}/(2d)\).
Let \(Y\) be the period lattice, an \(\mathcal O_E\)-module free
of finite rank \(\rho\) over \(\mathbb Z\). Write
\(B:Y\times Y\to\mathbb R\) for the positive definite symmetric
form supplied by the valuation of the period pairing. It is balanced
in the sense that \(B(sy,z)=B(y,\bar s z)\).
With the source's normalization \(Q(y)=B(y,y)/2\). Let
\(\mathcal S\) be the real vector space of balanced symmetric forms.
Let \(\phi(B)\) be the infimum of \(\lambda(B)\) over real linear
functionals on \(\mathcal S\) integral on the integral forms and
strictly positive on every nonzero positive semidefinite form.
The chart identifies this infimum with \(\phi_T\).
Such functionals exist: for an integral basis \(v_i\) of \(Y\) and
the host element \(w_0\in\mathcal O_E\) with \(w_0+\bar w_0=1\),
\(\lambda(C)=\sum_i C(w_0v_i,v_i)=\tfrac12\sum_i C(v_i,v_i)\)
has both properties. If the last sum vanishes for a positive
semidefinite \(C\), the whole basis lies in its radical.

There is a \(\mathbb Z\)-basis \(b_i\) of \(Y\) satisfying
\begin{equation}
  Q(b_i)\le R\rho^3\phi(B),\qquad
  |B(b_i,z)|\le 2\sqrt{\rho}\,Q(z)\quad(z\in Y).
  \label{eq:period-short-basis}
\end{equation}
The rank-zero case is empty. Assume \(\rho>0\).

\subsection{The dual minimum and the factor from complex multiplication}

Fix an admissible \(\lambda\), a basis \(v_j\) of \(Y\), and its
real dual coordinates \(v_j^*\). There are \(y_s,z_s\in Y\) with
\(\lambda(C)=\sum_s C(y_s,z_s)\) for every \(C\in\mathcal S\).
Indeed quotient \(Y\otimes_{\mathbb Z}Y\) by the common integral
annihilator of all the forms in \(\mathcal S\). The quotient is
torsion-free and finite free. The defining symmetry and balancing
relations have integral coefficients, so the real dual of this quotient
is precisely \(\mathcal S\): the common annihilator is the saturation
of the integral span of those relations. Integral forms are precisely
the integral dual. A functional integral on that dual is thus in the
quotient lattice, where its representative is a finite integral sum of
pure tensors. Absorb integral coefficients in the first factor.

In \(Y_{\mathbb C}\) let \(p_\sigma\) denote the idempotent of
\(E\otimes_{\mathbb Q}\mathbb C=\prod_\sigma\mathbb C\) corresponding
to \(\sigma\); write \(\eta_\sigma(y)=\eta_{\mathbb C}(p_\sigma y)\)
for a real covector \(\eta\). Define
\[
 C_\eta(y,z)=\sum_\sigma
       \operatorname{Re}(\eta_\sigma(y)\overline{\eta_\sigma(z)}),
       \qquad L[\eta]=\lambda(C_\eta).
\]
The form \(C_\eta\) is balanced and positive semidefinite, and
is nonzero when \(\eta\ne0\), by \(\sum_\sigma p_\sigma=1\).
Thus \(L\) is a positive definite quadratic form on the real dual.
In the chosen coordinates let \(L,G\) denote the matrices of this
quadratic form and of \(B\), respectively. Then
\(\operatorname{tr}(LG)=\lambda(B)\). Indeed polarize the
definition of \(C_\eta\), multiply its value on \(v_i^*,v_j^*\)
by \(B(v_i,v_j)\), and sum over \(i,j\).
The result at \(y,z\) before applying \(\lambda\) is
\(\sum_\sigma \operatorname{Re} B_{\mathbb C}
(p_\sigma y,\overline{p_\sigma z})=B(y,z)\).
The last equality follows from balancing, which pairs only
the \(\sigma\) and \(\bar\sigma\) components.

For an integral covector \(0\ne\eta\in Y^*\), let \(t_y\in E\)
represent the rational linear functional
\(c\mapsto\eta(cy)\) under the trace pairing. Then
\(t_y\in\mathfrak D_E^{-1}\) because this pairing with
\(\mathcal O_E\) is integral. In the idempotent coordinates
\(\eta_\sigma(y)=\sigma(t_y)\), by
\(\eta(cy)=\sum_\sigma \sigma(c)\sigma(t_y)\).
Some \(t_y\ne0\), and all of its embeddings are nonzero.
Set
\[
 t=\sum_s(t_{y_s}\overline{t_{z_s}}+
                     \overline{t_{y_s}}t_{z_s})
            \in\mathfrak D_E^{-2}.
\]
For each \(\sigma\), \(\sigma(t)\) is twice the value of \(\lambda\)
on \((y,z)\mapsto\operatorname{Re}
(\eta_\sigma(y)\overline{\eta_\sigma(z)})\), which is a nonzero
positive semidefinite balanced form. Thus every \(\sigma(t)>0\).
As \((t)\mathfrak D_E^2\) is a nonzero integral ideal,
\(\prod_\sigma \sigma(t)=|N_{E/\mathbb Q}t|\ge\Delta^{-2}\).
The arithmetic--geometric mean inequality gives
\begin{equation}
       L[\eta]=\tfrac12\sum_\sigma\sigma(t)
          \ge \tfrac d2\Delta^{-2/d}=:\mu^2 .
      \label{eq:period-dual-min}
\end{equation}

Take nonzero \(\theta\in\mathcal O_E\) with \(\bar\theta=-\theta\),
by subtracting the conjugate of a nonreal integral element.
Its action matrix \(K\) on \(Y\) is invertible over \(\mathbb R\).
Indeed \(Y_{\mathbb Q}\) is an \(E\)-module. Balancing gives
\(K^tG=-GK\); the polarized expression for \(L\), with
\(\overline{\sigma(\theta)}=-\sigma(\theta)\), gives \(KL=-LK^t\).
Write \(L=S^2\) with \(S\) positive definite symmetric, and put
\(D=SGS,\ J=S^{-1}KS\). Then \(J^t=-J\) and \(JD=DJ\).
Every eigenspace of \(D\) is preserved by the invertible
skew-symmetric \(J\); in particular its top eigenspace has even
dimension. As \(D\) is positive definite, its largest eigenvalue
is at most \(\tfrac12\operatorname{tr}D\).
For each column \(x\in\mathbb R^\rho\) this gives
\[
       x^tGx\le\tfrac12\lambda(B)\,\|S^{-1}x\|^2 .
\]

The Euclidean successive-minima transference inequality
\(\lambda_1(\Lambda)\lambda_\rho(\Lambda^*)\le\rho\)
applied to the lattice \(\Lambda=S\mathbb Z^\rho\) and its
reciprocal \(S^{-1}\mathbb Z^\rho\) now gives \(\rho\) independent
\(c_i\in\mathbb Z^\rho\) with
\(\|S^{-1}c_i\|\le\rho/\mu\), by \eqref{eq:period-dual-min}.
Discrete full lattices attain their successive minima.
In the \(B\)-norm, their associated lattice vectors therefore have
\[
    Q(c_i)\le \frac14\lambda(B)\,\rho^2/\mu^2
            =R\rho^2\lambda(B).
\]
The least possible maximum of \(Q\) on \(\rho\) independent
lattice vectors is attained, since every bounded \(B\)-ball in
\(Y\) is finite. The same minimum is thus at most
\(R\rho^2\phi(B)\) after taking the infimum over \(\lambda\).
\footnote{\href{https://doi.org/10.1007/BF01445125}{W. Banaszczyk,
New bounds in some transference theorems in the geometry of numbers,
Mathematische Annalen 296 (1993), 625--635}.
The archive gives a Gaussian Poisson proof of the transference input
in \href{https://github.com/ChristianKudeba/alpoge-rh-verification/blob/alpoge-source-7cbf7a0c/src/Formalism/PeriodIntegrals/OrbitalBoundsConclusion/TransferenceMinima.lean}{its transference theorem}.}

\subsection{Converting independent vectors into the required basis}

Take successive shortest vectors \(u_i\) outside the subspaces
\(V_{i-1}=\langle u_1,\ldots,u_{i-1}\rangle_{\mathbb R}\), and put
\(V_i=V_{i-1}+\mathbb R u_i\). Such vectors exist by discreteness.
The preceding independent tuple shows that each
\(\|u_i\|_B^2/2\le R\rho^2\phi(B)\), and the norms are nondecreasing.
The quotient \((Y\cap V_i)/(Y\cap V_{i-1})\) is free cyclic:
its denominator is saturated, and the ranks of the numerator and
denominator are \(i,i-1\). Choose a lift \(b_i\) of the positive
generator with smallest norm. By induction these lifts form an
integral basis on every step of the flag.

Let \(u_j^*\) be the Gram--Schmidt vectors for the \(u_j\),
with the same \(B\)-metric. The perpendicular component of \(b_i\)
along \(u_i^*\) is a positive multiple at most one of \(u_i^*\),
since the nonzero class of \(u_i\) is an integral multiple of the
generator. Subtracting integer multiples of
\(u_{i-1},\ldots,u_1\) successively reduces the earlier
Gram--Schmidt coordinates in the same lift to absolute value
at most \(1/2\). The minimum norm choice and orthogonality imply
\[
 \|b_i\|_B^2\le\|u_i^*\|_B^2+
       \tfrac14\sum_{j<i}\|u_j^*\|_B^2
       \le \rho\|u_i\|_B^2 .
\]
This proves the first inequality of \eqref{eq:period-short-basis}.
For \(z\in Y\cap V_{i-1}\), the lifts \(b_i+z,b_i-z\) have
norm at least \(\|b_i\|_B\). Expanding the two squares gives
\(2|B(b_i,z)|\le\|z\|_B^2\). For \(z\notin V_{i-1}\), minimality
gives \(\|u_i\|_B\le\|z\|_B\), so Cauchy--Schwarz gives
\(|B(b_i,z)|\le\sqrt{\rho}\,\|z\|_B^2\).
These bounds imply the second inequality in both cases.

In applying \eqref{eq:period-short-basis} to
Section~\ref{sec:cusp-jets}, the geometric inputs are the period
pairing with diagonal \(2Q\), identification of its infimum with
the frame slope \(\phi_T\), and the toric rank inequality
\(\rho\le d\,q_-\), where \(q_-\) denotes the negative signature
rank. The two numerical basis parameters are therefore exactly
\(R\rho^3\phi_T\) and \(2\sqrt{\rho}\).
\footnote{\href{https://github.com/ChristianKudeba/alpoge-rh-verification/blob/alpoge-source-7cbf7a0c/src/Formalism/SharpRestriction/Divisibility/GeneratorPowers/PhiCTie.lean}{Integral functionals on balanced forms},
\href{https://github.com/ChristianKudeba/alpoge-rh-verification/blob/alpoge-source-7cbf7a0c/src/Formalism/SharpRestriction/Divisibility/GeneratorPowers/PhiCMinI.lean}{the codifferent lower bound},
\href{https://github.com/ChristianKudeba/alpoge-rh-verification/blob/alpoge-source-7cbf7a0c/src/Formalism/SharpRestriction/Divisibility/GeneratorPowers/PhiCMinK.lean}{the trace identity, CM symmetry and transfer},
\href{https://github.com/ChristianKudeba/alpoge-rh-verification/blob/alpoge-source-7cbf7a0c/src/Formalism/SharpRestriction/Divisibility/GeneratorPowers/PhiCMinima.lean}{minimization over the cone of forms},
\href{https://github.com/ChristianKudeba/alpoge-rh-verification/blob/alpoge-source-7cbf7a0c/src/ReciprocityTheory/ReductiveGroups/ArithmeticGeometry/FirstCohomology/ReducedBasis.lean}{basis reduction and the mixed estimate},
\href{https://github.com/ChristianKudeba/alpoge-rh-verification/blob/alpoge-source-7cbf7a0c/src/Development/AbelianVarietyTheory/Proofs/DivisorProofs/ExistsShortPeriodLatticeBasisBalancedSlope.lean}{specialization to the period lattice of the chart}.}
